% =====================================================================
%  Partie 6 -- Cadre réglementaire
% =====================================================================
\part{Cadre réglementaire}
\chapter{Historique}
\chapter{La RE 2020}
\label{chap:re2020}

\section{Introduction}

\subsection*{Objectifs}
La France a adopté un objectif national de neutralité carbone en 2050, instauré par la loi Énergie-Climat (2019), renforcée par la Stratégie nationale bas carbone révisée (2020) et la Programmation pluriannuelle de l'énergie (2020).

Le secteur du bâtiment représente 45\,\% de la consommation énergétique nationale\footnote{ADEME : \texttt{https://www.ademe.fr/les-defis-de-la-transition/batiments} ; Batizoom (ADEME) : \texttt{https://batizoom.ademe.fr}} et 19\,\% des émissions de gaz à effet de serre (associées à l'utilisation de l'énergie)\footnote{ADEME, \emph{op. cit.} ; Xpair : \texttt{https://conseils.xpair.com/analyse\_reglementation/bilan-programme-obec-batiment-energie-carbone.htm}}.

\subsection*{RE et RT}
La France s'est dotée d'une nouvelle réglementation environnementale en 2022, qui prend progressivement la suite de la RT 2012 : la RE 2020.
Nous passons d'une Réglementation Thermique (RT) à une Réglementation Environnementale (RE).

\subsection*{RE 2020}
Elle poursuit des objectifs d'amélioration de la performance énergétique des bâtiments neufs, de réduction de leur impact sur le climat (prise en compte des émissions de gaz à effet de serre sur l'ensemble du cycle de vie des bâtiments) et de leur adaptation aux conditions climatiques futures (renforcement du confort d'été)\footnote{Guide de la RE2020 : \texttt{https://rt-re-batiment.developpement-durable.gouv.fr/IMG/pdf/guide\_re\_2020\_16mai2025.pdf}}.

Elle amène à une amélioration de la conception bioclimatique des bâtiments, elle renforce la performance de l'enveloppe du bâti, elle favorise le recours aux énergies renouvelables et peu carbonées et aux matériaux ayant une faible empreinte carbone.

\subsection*{Catégories et indicateurs}
La RE 2020 s'articule autour de trois catégories (carbone, énergie, confort d'été) et de six indicateurs (tableau \ref{tab:re2020-indicateurs}).

\begin{table}[htbp]
\centering
\caption{Les six indicateurs de la RE 2020}
\label{tab:re2020-indicateurs}
\small
\begin{tabular}{|p{1.6cm}|p{2.6cm}|p{3.8cm}|p{6cm}|}
\hline
\textbf{Catégorie} & \textbf{Indicateur} & \textbf{Définition} & \textbf{Description} \\
\hline
Carbone & $I_c$ construction (kgeqCO$_2$/m$^2$) & Impact sur le changement climatique associé aux « composants » + « chantier » & Introduction de la méthode d'analyse de cycle de vie pour l'évaluation des émissions de gaz à effet de serre des produits de construction et équipements et de leur mise en œuvre : « composants » et « chantier ». \\
\hline
Carbone & $I_c$ énergie (kgeqCO$_2$/m$^2$) & Impact sur le changement climatique associé aux consommations d'énergie primaire & Introduction de la méthode d'analyse de cycle de vie pour l'évaluation des émissions de gaz à effet de serre des énergies consommées pendant le fonctionnement du bâtiment, soit 50 ans. \\
\hline
Énergie & $C_{ep}$ (kWhep/(m$^2$.an)) & Consommation d'énergie primaire totale & Évaluation des consommations d'énergie renouvelable et non renouvelable pour : chauffage, refroidissement, ECS, éclairage, ventilation, éclairage et ventilation des parkings, éclairage des circulations en collectif, électricité des ascenseurs. \\
\hline
Énergie & $C_{ep,nr}$ (kWhep/(m$^2$.an)) & Consommation d'énergie primaire non renouvelable & Idem, pour la seule part non renouvelable. \\
\hline
Énergie & $B_{bio}$ (points) & Besoins bioclimatiques & Évaluation des besoins de chaud, de froid (que le bâtiment soit climatisé ou pas) et d'éclairage. \\
\hline
Confort d'été & DH (°C.h) & Degrés-heures d'inconfort : niveau d'inconfort perçu par les occupants sur l'ensemble de la saison chaude & Évaluation des écarts entre la température du bâtiment et une température de confort (température adaptée en fonction des températures des jours précédents, elle varie entre 26 et 28\,°C). \\
\hline
\end{tabular}
\end{table}

\subsection*{Progressivité}
La RE 2020 prévoit plusieurs phases de renforcement des critères carbone et énergie, en 2025, 2028 puis 2031.

La RE 2020 est entrée en vigueur le 1\textsuperscript{er} janvier 2022 pour les maisons individuelles et les logements collectifs, puis elle a progressivement intégré différents usages.
À partir du 1\textsuperscript{er} mai 2026, quasiment tous les programmes de constructions neuves sont concernés.

\section{Carbone}

\subsection*{Analyse du cycle de vie (ACV)}
L'\textbf{analyse du cycle de vie} (\textbf{ACV}) est une méthode d'évaluation permettant de réaliser un bilan environnemental \textbf{multicritère} d'un système (produit, service ou procédé) sur \textbf{l'ensemble de son cycle de vie} : extraction des matières premières, fabrication, transport, distribution, utilisation, fin de vie, valorisation.
% [Figure : cycle de vie]

\subsection*{L'unité fonctionnelle}
La norme ISO 14044 définit l'unité fonctionnelle comme la « performance quantifiée d'un système de produits, destinée à être utilisée comme unité de référence dans une ACV ».

En ACV, on ne compare pas des objets : on compare la même fonction rendue par des objets ou systèmes différents.
L'unité fonctionnelle est la traduction mesurable de cette fonction.

\subsection*{FDES}
Une Fiche de Déclaration Environnementale et Sanitaire (FDES) est une déclaration environnementale normée. Elle contient les résultats de l'analyse de cycle de vie d'un produit ainsi que des informations sanitaires.
Les FDES sont accessibles sur la base INIES : \texttt{https://base-inies.fr/consultation/recherche}.

\subsection*{$I_c$ projet}
L'impact carbone d'un projet se décompose comme indiqué au tableau \ref{tab:ic-projet} (exemple pour des bureaux).

\begin{table}[htbp]
\centering
\caption{Décomposition de l'$I_c$ projet (exemple : bureaux)}
\label{tab:ic-projet}
\small
\begin{tabular}{|p{3cm}|l|p{6cm}|c|}
\hline
\textbf{Regroupement} & \textbf{Composante} & \textbf{Contenu} & \textbf{Part} \\
\hline
$I_c$ construction & $I_c$ composant & Produits et équipements de construction & 65\,\% \\
(réglementaire) & $I_c$ chantier & Consommation d'énergie et d'eau sur le site, élimination et traitement des déchets & 3\,\% \\
\hline
$I_c$ énergie (réglementaire) & $I_c$ énergie & Consommation d'énergie & 25\,\% \\
\hline
 & $I_c$ eau & Consommation d'eau & 7\,\% \\
\hline
 & $I_c$ parcelle & Utilisation de la parcelle, arrosage & -- \\
\hline
\end{tabular}
\end{table}
$I_c$ bâtiment = $I_c$ construction + $I_c$ énergie + $I_c$ eau ; $I_c$ projet = $I_c$ bâtiment + $I_c$ parcelle.

\subsection*{$I_c$ construction}
L'$I_c$ construction (kgeqCO$_2$/m$^2$) représente l'impact de la construction sur le changement climatique : matériaux et équipements + chantier. Il doit respecter :
\begin{equation}
  I_{c,\mathrm{construction}} \leq I_{c,\mathrm{construction},\mathrm{maxmoyen}} \times \left(1 + Mi_{\mathrm{combles}} + Mi_{\mathrm{surf}}\right) + Mi_{\mathrm{infra}} + Mi_{\mathrm{vrd}} + Mi_{\textrm{géo}} + Mi_{\mathrm{ded}}
\end{equation}
Le seuil moyen est évolutif. Les modulations sont les suivantes :
\begin{itemize}
  \item \textbf{Combles} : permet d'ajuster les exigences pour prendre en compte les surfaces construites non habitables ($h < 1{,}80$\,m).
  \item \textbf{Surface} : permet d'équilibrer les efforts entre petits et grands bâtiments (moyenne : $-30$\,kgCO$_2$e/m$^2$).
  \item \textbf{Infrastructure} : permet d'ajuster les exigences pour prendre en compte les contraintes nécessitant d'importantes fondations ou des espaces souterrains (moyenne : 45\,kgCO$_2$e/m$^2$).
  \item \textbf{VRD} : permet d'ajuster les exigences pour tenir compte des impacts élevés liés à la présence de réseaux d'alimentation du bâtiment ou d'aires de stationnement qui desservent le bâtiment (moyenne : 6\,kgCO$_2$e/m$^2$).
  \item \textbf{Géographie} : permet d'ajuster les exigences pour tenir compte de la mise en œuvre de moyens supplémentaires pour assurer le confort estival du pourtour et de l'arrière-pays méditerranéens.
  \item \textbf{DED} : permet d'ajuster les exigences pour prendre en compte le manque de disponibilité de données environnementales spécifiques (moyenne : 12\,kgCO$_2$e/m$^2$).
\end{itemize}
Source : Observatoire de la RE2020 -- Hub bas carbone -- IFPEB.

\begin{table}[htbp]
\centering
\caption{Seuils $I_{c,\mathrm{construction},\mathrm{maxmoyen}}$ (kgCO$_2$e/m$^2$)}
\label{tab:ic-construction-seuils}
\begin{tabular}{|l|c|c|c|c|}
\hline
\textbf{Usage} & \textbf{2022} & \textbf{2025} & \textbf{2028} & \textbf{2031} \\
\hline
Logement & 740 & 650 & 580 & 490 \\
Bureau & 980 & 810 & 710 & 600 \\
Enseignement primaire et secondaire & 900 & 770 & 680 & 590 \\
\hline
\end{tabular}
\end{table}

\subsection*{$I_c$ énergie}
L'$I_c$ énergie (kgeqCO$_2$/m$^2$) représente l'impact de l'énergie sur le changement climatique (consommation d'énergie et d'eau). Il doit respecter :
\begin{equation}
  I_{c,\textrm{énergie}} \leq I_{c,\textrm{énergie},\mathrm{maxmoyen}} \times \left(1 + Mi_{\mathrm{combles}} + Mi_{\mathrm{surf\_moy}} + Mi_{\mathrm{surf\_tot}} + Mi_{\mathrm{cat}} + Mi_{\textrm{géo}}\right)
\end{equation}
Les seuils moyens sont évolutifs et donnés par le gouvernement par typologie de programme. Les modulations sont les suivantes :
\begin{itemize}
  \item \textbf{Combles} : permet d'ajuster les exigences pour prendre en compte les surfaces construites non habitables ($h < 1{,}80$\,m).
  \item \textbf{Surface} : permet d'équilibrer les efforts entre petits et grands bâtiments.
  \item \textbf{Surface totale} : modulation spécifique pour les logements collectifs pour équilibrer les efforts liés aux parties collectives entre petits et grands collectifs.
  \item \textbf{Contraintes} : permet de compenser les contraintes extérieures qui limitent les possibilités de ventilation naturelle du bâtiment par ouverture des fenêtres.
  \item \textbf{Géographie} : permet d'ajuster les exigences pour tenir compte de la mise en œuvre de moyens supplémentaires pour assurer le confort estival du pourtour et de l'arrière-pays méditerranéens.
\end{itemize}

\begin{table}[htbp]
\centering
\caption{Seuils $I_{c,\textrm{énergie},\mathrm{maxmoyen}}$ (kgCO$_2$e/m$^2$)}
\label{tab:ic-energie-seuils}
\begin{tabular}{|l|c|c|c|}
\hline
\textbf{Usage} & \textbf{2022} & \textbf{2025} & \textbf{2028} \\
\hline
Logement & 560 & 320 & 260 \\
Bureau & 280 & 200 & -- \\
Enseignement primaire et secondaire & 240 & 200 & 140 \\
\hline
\end{tabular}
\end{table}

\subsection*{Un outil en ligne pour l'ACV et l'$I_c$}
Décarbone+ (AIA Environnement) : \texttt{https://www.decarbone-plus.com/}

\section{Énergie}

\subsection*{Énergie primaire}
L'énergie primaire est une forme d'énergie disponible dans la nature avant toute transformation.

En France, la consommation d'énergie primaire s'élève à 2\,630\,TWh en 2024 (donnée corrigée des variations climatiques)\footnote{\texttt{https://www.statistiques.developpement-durable.gouv.fr/}}, répartie en :
\begin{itemize}
  \item 41\,\% de nucléaire ;
  \item 28\,\% de pétrole ;
  \item 12\,\% de gaz naturel ;
  \item 17\,\% d'énergies renouvelables et déchets (dont 1\,\% de déchets non renouvelables) ;
  \item 2\,\% de charbon.
\end{itemize}
% [Figure : évolution de la consommation d'énergie primaire 1990-2024]

\subsection*{Consommation d'énergie primaire}
La consommation d'énergie primaire $C_{ep}$ (kWhep/(m$^2$.an)) doit respecter :
\begin{equation}
  C_{ep} \leq C_{ep,\mathrm{maxmoyen}} \times \left(1 + Mi_{\mathrm{combles}} + Mi_{\mathrm{surf\_moy}} + Mi_{\mathrm{surf\_tot}} + Mi_{\mathrm{cat}} + Mi_{\textrm{géo}}\right)
\end{equation}
Les modulations sont les mêmes que pour l'$I_c$ énergie (combles, surface, surface totale, contraintes, géographie).

\subsection*{Consommation d'énergie primaire non renouvelable}
De même, la consommation d'énergie primaire non renouvelable $C_{ep,nr}$ (kWhep/(m$^2$.an)) doit respecter :
\begin{equation}
  C_{ep,nr} \leq C_{ep,nr,\mathrm{maxmoyen}} \times \left(1 + Mi_{\mathrm{combles}} + Mi_{\mathrm{surf\_moy}} + Mi_{\mathrm{surf\_tot}} + Mi_{\mathrm{cat}} + Mi_{\textrm{géo}}\right)
\end{equation}

\begin{table}[htbp]
\centering
\caption{Seuils de consommation d'énergie primaire (kWhep/(m$^2$.an))}
\label{tab:cep-seuils}
\begin{tabular}{|l|c|c|}
\hline
\textbf{Usage} & \textbf{$C_{ep,\mathrm{maxmoyen}}$} & \textbf{$C_{ep,nr,\mathrm{maxmoyen}}$} \\
\hline
Logement & 85 & 70 \\
Bureau & 85 & 75 \\
Enseignement primaire et secondaire & 72 & 65 \\
\hline
\end{tabular}
\end{table}

\subsection*{Besoin bioclimatique conventionnel}
Le calcul du $B_{bio}$ RE2020 comprend les besoins en chauffage, en refroidissement et en éclairage du bâtiment.
Le $B_{bio}$ s'intéresse uniquement au bâti et non aux équipements.

Le $B_{bio}$ s'intéresse à :
\begin{itemize}
  \item l'orientation ;
  \item l'isolation ;
  \item la compacité ;
  \item les surfaces vitrées ;
  \item la protection solaire ;
  \item l'inertie ;
  \item la conception globale de l'enveloppe.
\end{itemize}

Le $B_{bio}$ (en points) doit respecter :
\begin{equation}
  B_{bio} \leq B_{bio,\mathrm{maxmoyen}} \times \left(1 + Mi_{\mathrm{combles}} + Mi_{\mathrm{surf\_moy}} + Mi_{\mathrm{surf\_tot}} + Mi_{\mathrm{bruit}} + Mi_{\textrm{géo}}\right)
\end{equation}
Les modulations sont : combles, surface, surface totale, géographie (comme précédemment), et \textbf{bruit} : coefficient de modulation selon l'exposition au bruit des infrastructures de transport à proximité du bâtiment.

\section{Confort d'été}

\subsection*{Degrés-heures d'inconfort}
Le calcul des degrés-heures d'inconfort (DH) suit les étapes suivantes :
\begin{enumerate}
  \item météo avec scénario caniculaire, en fonction de la zone climatique ;
  \item calcul en mode confort d'été (simulation avec un logiciel spécialisé) ;
  \item somme des degrés-heures ;
  \item comparaison du DH du projet au DH$_{\max}$.
\end{enumerate}

Chaque heure, on compte l'écart entre la température du bâtiment et la température limite de confort (lorsque la première dépasse la seconde). Par exemple, pour une journée type d'été en occupation, avec un seuil à 27\,°C (température extérieure moyenne des jours passés haute), le cumul de la journée atteint 20,0\,°C.h.
% [Figure : journée type d'été en occupation]

\subsection*{Seuils}
\begin{itemize}
  \item \textbf{Catégorie 1} : situation à moindre contrainte extérieure, ainsi que les bâtiments non climatisés (sauf parties de bâtiments climatisées en zones H2d et H3).
  \item \textbf{Catégorie 1 climatisée} : pour les zones H2d et H3.
  \item \textbf{Catégorie 2} : situation de contrainte forte. Les bâtiments doivent être localisés dans une région « chaude » à faible altitude, climatisés et exposés à des nuisances sonores.
  \item \textbf{Catégorie 3} : situation de contrainte très forte, uniquement pour les bureaux climatisés soumis à des règles interdisant l'ouverture des fenêtres et les immeubles de grande hauteur.
\end{itemize}

\begin{table}[htbp]
\centering
\caption{Seuils de degrés-heures d'inconfort DH$_{\max}$ (°C.h)}
\label{tab:dh-seuils}
\small
\begin{tabular}{|p{4.2cm}|c|c|c|c|}
\hline
\textbf{Usage} & \textbf{Cat. 1} & \textbf{Cat. 1 clim.} & \textbf{Cat. 2} & \textbf{Cat. 3} \\
\hline
Logement, $S_{moy} \leq 20$\,m$^2$ & 1250 & 1600 & 2600 & -- \\
Logement, $20 < S_{moy} \leq 60$\,m$^2$ & 1250 & $1700 - 5\times S_{moy}$ & $2850 - 12{,}5\times S_{moy}$ & -- \\
Logement, $S_{moy} > 60$\,m$^2$ & 1250 & 1400 & 2100 & -- \\
\hline
Bureau & 1150 & 2400 & 2600 & pas de seuil \\
\hline
Enseignement primaire et secondaire & 900 & 1800 & 2200 & -- \\
\hline
\end{tabular}
\end{table}
